\subsection{子模的有序集}

在本小节中，A 与 B 是 k-代数，M 是一个左 A-模，V 是一个左 B-模。我们用 \(\mathrm{D}(\mathrm{M})\)（相应地，\(\mathrm{D}(\mathrm{V})\)）表示 M（相应地，V）的按包含关系排序的子模之集。我们设给定了有序集的同构 \(\varphi : \mathrm{D}(\mathrm{V}) \to \mathrm{D}(\mathrm{M})\)。

由森田定理（VIII, 第 103 页，定理 2），在下列情形我们得到这样一个同构：P 是一个可逆 \((\mathrm{A},\mathrm{B})_k\)-双模，M 是 A-模 \(\mathrm{P} \otimes_{\mathrm{B}} \mathrm{V}\)，并且对 V 的每个子模 W，\(\varphi(\mathrm{W})\) 是 \(\mathrm{P} \otimes_{\mathrm{B}} \mathrm{W}\) 在 M 中的典范像。

模 M 的某些性质，或其子模的某些性质，可以用有序集 D(M) 来表达：它们列于表 I 与表 II 中。

模 M 是子模族 \((\mathrm{M}_{i})_{i\in\mathrm{I}}\) 的直和当且仅当我们有 \(M=\sum_{i\in I}M_{i}\) 与 \(M_{i}\cap\sum_{j\neq i}M_{j}=0\)（对每个 \(i\in I\)）。这一注记以及对表 I 的考察给出下面的结果。

命题 7. —— a) 我们有 \(\varphi(0)=0\) 及 \(\varphi(V)=M\)。

\begin{enumerate}
\def\labelenumi{\alph{enumi})}
\setcounter{enumi}{1}
\tightlist
\item
  设 \((\mathrm{V}_i)_{i\in \mathrm{I}}\) 是 V 的一个子模族。我们有
\end{enumerate}

\[
\varphi \left(\sum_ {i \in \mathrm{I}} \mathrm{V} _ {i}\right) = \sum_ {i \in \mathrm{I}} \varphi (\mathrm{V} _ {i}), \quad \varphi \left(\bigcap_ {i \in \mathrm{I}} \mathrm{V} _ {i}\right) = \bigcap_ {i \in \mathrm{I}} \varphi (\mathrm{V} _ {i}).
\]

M 的子模

有序集 D(M)

零子模

D(M) 的最小元

子模 M

D(M) 的最大元

\(\bigcap_{i\in I} M_i\)

下确界 \(\inf_{i\in I} M_i\)

\(\sum_{i\in I} M_i\)

上确界 \(\sup_{i\in I} M_i\)

互补子模

\(\inf(M', M'') = 0, \sup(M', M'') = M\)

M 的单子模

D(M) --- \{0\} 的极小元

M 的极大子模

D(M) --- \{M\} 的极大元

M 的底座 \(\mathscr{S}(M)\)

D(M) --- \{0\} 的极小元之集在 D(M) 中的上确界

\emph{M 的根 \(\Re(M)\)}（VIII, 第 151 页）

D(M) --- \{M\} 的极大元之集在 D(M) 中的下确界

模 M 的性质

D(M) 的性质

M 是诺特模。

有序集 D(M) 是诺特的（《集合论》III, §6, No.~5, 第 190 页）。

M 是阿廷模。

按 ⊃ 排序的集合 D(M) 是诺特的。

M 是不可分解的。

我们有 M ≠ 0，并且不存在 D(M) 的两个非零元素 M\textquotesingle{} 与 M\textquotesingle\textquotesingle{} 满足 inf(M\textquotesingle, M\textquotesingle\textquotesingle) = 0, sup(M\textquotesingle, M\textquotesingle\textquotesingle) = M。

M 是有限生成的。

对 D(M) 中每个以 M 为上界的族 (Mi) i∈I，存在 I 的有限子集 J，使 M = supj∈I Mj。

M 是单模。

Card(D(M)) = 2。

M 是半单的。

M 是 D(M) = \{0\} 的极小元之集在 D(M) 中的上确界。

表 II.

\begin{enumerate}
\def\labelenumi{\alph{enumi})}
\setcounter{enumi}{2}
\tightlist
\item
  B-模 V 是子模族 \((\mathrm{V}_{i})_{i\in\mathrm{I}}\) 的直和当且仅当 M 是子模族 \((\varphi(\mathrm{V}_{i}))_{i\in\mathrm{I}}\) 的直和。
\end{enumerate}

设 \(V'\) 与 \(V''\) 是 \(V\) 的子模，使得 \(V'\) 含于 \(V''\)；令 \(M' = \varphi(V')\)，\(M'' = \varphi(V'')\)，于是 \(M''\) 包含 \(M'\)。用 \([V', V'']\) 表示 \(D(V)\) 中由子模 \(W\)（属于 \(V\)，满足 \(V' \subset W \subset V''\)）组成的区间，并同样地定义区间 \([M', M'']\)（在 \(D(M)\) 中）。映射 \(W \mapsto W/V'\) 是从 \([V', V'']\) 到 \(D(V''/V')\) 的有序集同构；我们同样地定义从 \([M', M'']\) 到 \(D(M''/M')\) 的有序集同构。由于 \(\varphi\) 把区间 \([V', V'']\) 映到 \([M', M'']\)，它定义了一个有序集同构 \(\varphi\)，从 \(D(V''/V')\) 到 \(D(M''/M')\)。我们由此以及表 I 与表 II 推出下面的命题。

命题 8. —— a) 设 V' 与 V'\,' 是 V 的子模，使得 V'\,' 包含 V'。B-模 V'\,'/V' 是单模当且仅当 A-模 \(\varphi(\mathrm{V}^{\prime\prime})/\varphi(\mathrm{V}^{\prime})\) 是单模。

\begin{enumerate}
\def\labelenumi{\alph{enumi})}
\setcounter{enumi}{1}
\item
  若 \(V'\) 是 V 的单子模、极大子模或直因子，则 \(\varphi(V')\) 相应地是 M 的单子模、极大子模或直因子。
\item
  同构 \(\varphi\) 把底座 \(\mathcal{S}(\mathrm{V})\)（\(\mathrm{V}\) 的）变为底座 \(\mathcal{S}(\mathrm{M})\)（\(\mathrm{M}^*\) 的），把根（VIII, 第 151 页，定义 1）\(\Re (\mathrm{V})\)（\(\mathrm{V}\) 的）变为根 \(\Re (\mathrm{M})\)（\(\mathrm{M}_{*}\) 的）
\item
  设 \((\mathrm{V}_{i})_{0\leqslant i\leqslant n}\) 是 V 的子模的一个有限序列。它是 V 的若尔当–赫尔德列当且仅当 \((\varphi(\mathrm{V}_{i})_{0\leqslant i\leqslant n})\) 是 M 的若尔当–赫尔德列。
\end{enumerate}

引理 2. —— 设 H 与 H' 是 V 的子模，满足 \(H \cap H' = 0\)。B-模 H 与 H' 同构当且仅当 A-模 \(\varphi(\mathrm{H})\) 与 \(\varphi(\mathrm{H}')\) 同构。

我们把 \(H+H'\) 等同于乘积 \(H \times H'\)。从 H 到 \(H'\) 的一个同构的图像是 V 的一个子模 \(H''\)，它满足

\[
\mathrm{H} \cap \mathrm{H} ^ {\prime \prime} = \mathrm{H} ^ {\prime} \cap \mathrm{H} ^ {\prime \prime} = 0, \quad \mathrm{H} + \mathrm{H} ^ {\prime} = \mathrm{H} + \mathrm{H} ^ {\prime \prime} = \mathrm{H} ^ {\prime} + \mathrm{H} ^ {\prime \prime};\tag{17}
\]

反之，每个具有这些性质的子模是从 H 到 \(H'\) 的一个同构的图像。由命题 7，关系 \(H \cap H' = 0\) 等价于 \(\varphi(\mathrm{H}) \cap \varphi(\mathrm{H}') = 0\)，而关系式 (17) 等价于关系式

\[
\varphi (\mathrm{H}) \cap \varphi (\mathrm{H} ^ {\prime \prime}) = \varphi (\mathrm{H} ^ {\prime}) \cap \varphi (\mathrm{H} ^ {\prime \prime}) = 0,
\]

\[
\varphi (\mathrm{H}) + \varphi (\mathrm{H} ^ {\prime}) = \varphi (\mathrm{H}) + \varphi (\mathrm{H} ^ {\prime \prime}) = \varphi (\mathrm{H} ^ {\prime}) + \varphi (\mathrm{H} ^ {\prime \prime});
\]

引理由此得出。

命题 9. —— 设 S 是 V 的一个单子模，T 是 M 的单子模 \(\varphi(\mathrm{S})\)。若用 \(V_{S}\) 表示 V 的 S 型同型分量，用 \(M_{T}\) 表示 M 的 T 型同型分量，则我们有 \(\varphi(V_{\mathrm{S}})=M_{\mathrm{T}}\)。

V 的每个异于 S 的单子模 \(S'\) 满足 \(S' \cap S = 0\)。因此它同构于 S 当且仅当 \(\varphi(S')\) 同构于 T（引理 2）。而 \(V_{S}\) 是 V 的同构于 S 的单子模之和，\(M_{T}\) 是 M 的同构于 T 的单子模之和。于是命题 9 由命题 7 与命题 8 立即得出。

命题 10. —— a) B-模 V 是阿廷的（相应地，诺特的、不可分解的、单的、有限生成的）当且仅当 M 亦然。

\begin{enumerate}
\def\labelenumi{\alph{enumi})}
\setcounter{enumi}{1}
\item
  B-模 V 具有有限长度当且仅当 A-模 M 具有有限长度，且此时我们有 \(\mathrm{long}_{\mathrm{B}}(\mathrm{V}) = \mathrm{long}_{\mathrm{A}}(\mathrm{M})\)。
\item
  B-模 V 是半单的（相应地，同型的）当且仅当 A-模 M 是半单的（相应地，同型的）。若情形如此，则我们有 \(\text{long}_{\text{B}}(\text{V}) = \text{long}_{\text{A}}(\text{M})\)。
\end{enumerate}

断言 a) 由表 II 中所列的第二条性质得出。

断言 b) 由命题 8, d) 得出。

模 V 是半单的当且仅当它等于其底座 \(\mathcal{S}(V)\)；它是同型的当且仅当存在 V 的一个单子模 S，使得 \(V = V_{S}\)。于是断言 c) 由命题 7, c)、8, c) 与 9（VIII, 第 106 页与第 109 页）得出。

